% Generated draft: requires mathematical and source review.
\documentclass[11pt]{article}
\usepackage[margin=25mm]{geometry}
\usepackage{fontspec}
\setmainfont{Latin Modern Roman}
\usepackage{amsmath,amssymb,mathtools,unicode-math}
\setmathfont{Latin Modern Math}
\usepackage{enumitem,needspace,xurl,hyperref}
\setlength{\parindent}{0pt}
\setlength{\parskip}{5pt}
\setlength{\emergencystretch}{3em}
\newcommand{\N}{\mathbb{N}}
\newcommand{\Z}{\mathbb{Z}}
\newcommand{\Q}{\mathbb{Q}}
\newcommand{\R}{\mathbb{R}}
\newcommand{\C}{\mathbb{C}}
\newcommand{\F}{\mathbb{F}}
\newcommand{\A}{\mathbb{A}}
\newcommand{\PP}{\mathbb{P}}
\newcommand{\E}{\mathbb{E}}
\newcommand{\eps}{\varepsilon}
\newcommand{\dd}{\,\mathrm{d}}
\begin{document}
\section*{UnsolvedMath ID: 800555}
\textbf{Source problem code:} catalogue entry
\subsection{Definitions and mathematical statement}
Let \(G=(V,E)\) be a simple undirected graph. A rectangle representation assigns to each \(v\in V\) a closed axis-parallel rectangle
\[
R_v=[a_v,b_v]\times[c_v,d_v]\subset\R^2,
\qquad a_v<b_v,\quad c_v<d_v,
\]
with pairwise disjoint interiors. Here axis-parallel is the meaning of the source's term \emph{isothetic}, and the strict inequalities express its non-zero-thickness requirement.

Two rectangles are horizontally or vertically visible if a horizontal or vertical beam of unobstructed visibility of positive finite transverse width connects them. Such a beam is a rectangular corridor connecting facing sides of the two rectangles, with no intervening rectangle blocking the corridor. Thus visibility requires positive transverse width, rather than merely an unobstructed line segment. The source does not specify the treatment of touching rectangles, zero-length corridors, or contact between a beam boundary and another closed rectangle; these conventions remain unresolved and must be fixed before a fully precise geometric predicate can be given.

Subject to those unresolved boundary conventions, \(G\) is a \emph{rectangle visibility graph} if it admits such a representation satisfying, for every pair of distinct vertices \(u,v\in V\),
\[
\{u,v\}\in E
\quad\Longleftrightarrow\quad
R_u\text{ and }R_v\text{ are horizontally or vertically visible}.
\]
\textbf{Problem 29 (Ghosh and Goswami, 2010).} Characterize rectangle visibility graphs: give necessary and sufficient graph-theoretic conditions for a graph to admit a rectangle representation with exactly these visibility adjacencies. The source does not explicitly specify whether graphs must be finite.

The supplied discussion presents this class as a generalization of bar visibility graphs and reports that no characterization was known and that Shermer had shown recognition to be NP-complete. These are supplied historical claims, not independently verified present-day assessments. The local registry labels the problem open.
\subsection{Short English statement}
Characterize the graphs representable by closed axis-parallel rectangles with pairwise disjoint interiors, where two vertices are adjacent exactly when their rectangles see each other horizontally or vertically through an unobstructed beam of positive finite width. The supplied formulation leaves contact conventions and graph finiteness unspecified. This is a statement-only draft requiring human mathematical review.
\subsection{Sources}
\begin{itemize}
\item Local UnsolvedMath registry, permanent ID \(800555\), source problem code \texttt{catalogue entry}, title \emph{Visibility Graphs of Points, Segments and Polygons --- Problem 29}. The supplied record attributes the problem to Ghosh and Goswami, gives the year 2010, and labels its status \texttt{open}; this label is supplied metadata rather than an independently verified mathematical assessment.
\item Ghosh and Goswami, \emph{Unsolved Problems in Visibility Graphs of Points, Segments and Polygons} (2010), Problem 29, \url{https://arxiv.org/abs/1012.5187}. The supplied provenance records access on 2026-10-16 and extraction by \texttt{arxiv\_latex}; these are supplied provenance statements, not a claim of access or verification for this draft.
\item The source cites the following reference labels for rectangle visibility graphs: \texttt{bdhs-orvg-1997}, \texttt{dh-rvrbp-1997}, \texttt{o-opcv-1998}, \texttt{r-drvg-1997}, and \texttt{sw-rvgc-2003}. Their bibliographic entries are not included in the supplied record.
\item The source attributes NP-completeness of recognition to Shermer under reference label \texttt{s-orvsiii-1996}. Its bibliographic entry and precise theorem hypotheses are not included in the supplied record, and the attribution has not been independently verified here.
\end{itemize}
\end{document}
